\documentclass[12pt,a4paper]{article}
\usepackage[a4paper,top=15mm,bottom=15mm,left=16mm,right=16mm]{geometry}
\usepackage[T1]{fontenc}
\usepackage{amsmath}
\usepackage{amssymb}
\usepackage{enumitem}

\setlength{\parindent}{0pt}
\setlength{\parskip}{4pt}
\allowdisplaybreaks

\begin{document}

\begin{center}
{\Large\bfseries Recurring Decimals into Fractions}\\[2pt]
{\large Answer Key}
\end{center}

\medskip
\noindent\textbf{\large Part 1: Reading and Writing Recurring Decimals}

\medskip
\noindent\textbf{Q1.} First eight decimal places:
\begin{enumerate}[label=(\alph*),leftmargin=*]
\item $0.\dot{5}=0.55555555$
\item $0.\dot{1}\dot{8}=0.18181818$
\item $0.3\dot{7}=0.37777777$
\item $0.\dot{2}0\dot{6}=0.20620620$
\item $0.4\dot{5}\dot{1}=0.45151515$
\item $0.\dot{9}=0.99999999$
\end{enumerate}

\medskip
\noindent\textbf{Q2.} Dot notation:
\begin{enumerate}[label=(\alph*),leftmargin=*]
\item $0.8888\ldots=0.\dot{8}$
\item $0.636363\ldots=0.\dot{6}\dot{3}$
\item $0.24444\ldots=0.2\dot{4}$
\item $0.157157\ldots=0.\dot{1}5\dot{7}$
\item $0.91666\ldots=0.91\dot{6}$
\item $0.40505\ldots=0.4\dot{0}\dot{5}$
\end{enumerate}

\medskip
\noindent\textbf{Q3.}
\begin{enumerate}[label=(\alph*),leftmargin=*]
\item $0.4\dot{5}$ is bigger.
\item One example: $0.\dot{3}\dot{5}$ lies between $0.\dot{3}$ and $0.\dot{4}$.
\end{enumerate}

\medskip
\noindent\textbf{\large Part 2: Finding Fractions from Recurring Decimals}

\medskip
\noindent\textbf{Q4.} $0.\dot{7}=\dfrac{7}{9}$

\medskip
\noindent\textbf{Q5.} $0.\dot{8}=\dfrac{8}{9}$

\medskip
\noindent\textbf{Q6.}
\begin{enumerate}[label=(\alph*),leftmargin=*]
\item $\dfrac{2}{9}$
\item $\dfrac{5}{9}$
\item $\dfrac{2}{3}$
\item $\dfrac{1}{3}$
\end{enumerate}

\medskip
\noindent\textbf{\large Part 3: Longer Repeating Blocks}

\medskip
\noindent\textbf{Q7.} $0.\dot{7}\dot{2}=\dfrac{8}{11}$

\medskip
\noindent\textbf{Q8.}
\begin{enumerate}[label=(\alph*),leftmargin=*]
\item $\dfrac{5}{33}$
\item $\dfrac{2}{33}$
\item $\dfrac{6}{11}$
\item $\dfrac{10}{11}$
\item $\dfrac{4}{37}$
\end{enumerate}

\medskip
\noindent\textbf{Q9.}
\begin{enumerate}[label=(\alph*),leftmargin=*]
\item Multiply by $10^n$, where $n$ is the number of digits in the repeating block.
\item $0.\dot{1}\dot{2}=\dfrac{12}{99}=\dfrac{4}{33}$. Ben and Kayla are both correct; Ben has just simplified Kayla’s fraction.
\end{enumerate}

\medskip
\noindent\textbf{\large Part 4: When the Repeating Part Does Not Start Straight Away}

\medskip
\noindent\textbf{Q10.} $0.2\dot{7}=\dfrac{5}{18}$

\medskip
\noindent\textbf{Q11.}
\begin{enumerate}[label=(\alph*),leftmargin=*]
\item $\dfrac{13}{30}$
\item $\dfrac{7}{12}$
\item $\dfrac{1}{12}$
\item $\dfrac{41}{330}$
\end{enumerate}

\medskip
\noindent\textbf{Q12.}
\begin{enumerate}[label=(\alph*),leftmargin=*]
\item $\dfrac{25}{9}$
\item $\dfrac{16}{11}$
\item $\dfrac{35}{11}$
\end{enumerate}

\medskip
\noindent\textbf{Q13.}
\begin{enumerate}[label=(\alph*),leftmargin=*]
\item $\dfrac{8}{9}$
\item $\dfrac{7}{11}$
\item $\dfrac{23}{90}$
\item $\dfrac{1}{37}$
\item $\dfrac{13}{18}$
\item $\dfrac{4}{33}$
\item $\dfrac{5}{12}$
\item $\dfrac{29}{450}$
\end{enumerate}

\medskip
\noindent\textbf{\large Part 5: Challenges}

\medskip
\noindent\textbf{Q14.}
\begin{enumerate}[label=(\alph*),leftmargin=*]
\item $0.\dot{9}=1$
\item Since $10x=9.\dot{9}$, subtracting gives $9x=9$, so $x=1$. Also $3\times0.\dot{3}=1$. There is no “just below” $1$; between any number below $1$ and $1$ there is always another number.
\item No. For any $a<1$, we have $a<\dfrac{a+1}{2}<1$. So there is no largest number smaller than $1$.
\end{enumerate}

\medskip
\noindent\textbf{Q15.}
\begin{enumerate}[label=(\alph*),leftmargin=*]
\item
$\dfrac12=0.5,\quad
\dfrac13=0.\dot{3},\quad
\dfrac14=0.25,\quad
\dfrac15=0.2,\quad
\dfrac16=0.1\dot{6},$

$\dfrac17=0.\dot{1}4285\dot{7},\quad
\dfrac18=0.125,\quad
\dfrac19=0.\dot{1},\quad
\dfrac1{10}=0.1,$

$\dfrac1{11}=0.\dot{0}\dot{9},\quad
\dfrac1{12}=0.08\dot{3},\quad
\dfrac1{16}=0.0625.$

Terminating decimals:
$\dfrac12,\ \dfrac14,\ \dfrac15,\ \dfrac18,\ \dfrac1{10},\ \dfrac1{16}$.
\item The terminating denominators have only prime factors $2$ and $5$.
\item $\dfrac1{15}$ recurs, $\dfrac1{40}$ terminates, $\dfrac1{35}$ recurs, $\dfrac1{64}$ terminates.
\item A fraction in lowest terms terminates exactly when its denominator has no prime factors other than $2$ and $5$.
\end{enumerate}

\medskip
\noindent\textbf{Q16.}
\begin{enumerate}[label=(\alph*),leftmargin=*]
\item
$\dfrac27=0.\dot{2}8571\dot{4},$

$\dfrac37=0.\dot{4}2857\dot{1},$

$\dfrac47=0.\dot{5}7142\dot{8},$

$\dfrac57=0.\dot{7}1428\dot{5},$

$\dfrac67=0.\dot{8}5714\dot{2}.$
\item They all use the same six digits in the same cyclic order:
\[
1\to4\to2\to8\to5\to7\to1.
\]
\item $x=0.\dot{2}8571\dot{4}=\dfrac{285714}{999999}=\dfrac27$.
\end{enumerate}

\medskip
\noindent\textbf{Q17.}
\begin{enumerate}[label=(\alph*),leftmargin=*]
\item One valid example is $0.\dot{2}\dot{7}=\dfrac{3}{11}$.
\item $\dfrac56=0.8\dot{3}$.
\item Yes. If the repeating block has $n$ digits and represents the number $N$, then
\[
10^n x=N+x,
\]
so
\[
x=\frac{N}{10^n-1}.
\]
Non-repeating leading digits can first be shifted past by multiplying by a power of $10$. Hence every recurring decimal is rational. Since $\pi$ is irrational, no recurring decimal can be exactly equal to $\pi$.
\end{enumerate}

\end{document}